\section{Absolute value, squares, and intervals}\label{sec:magnitude}
This section puts the logic of the chapter to work on real numbers. Inequalities that involve squares are a common place to lose half of an answer, and the words ``and'' and ``or'' are what keep track of both halves.

The \emph{absolute value} $|x|$ of a real number $x$ is its distance from zero on the number line. Thus $|3|=3$, $|-3|=3$, and $|0|=0$. As a formula,
\[
|x|=\begin{cases}x&\text{if }x\ge0,\\-x&\text{if }x<0.\end{cases}
\]
The large brace lists two cases: use the first line when $x\ge0$ and the second line when $x<0$. In the second case the result is positive, because the negative of a negative number is positive; for instance, $|-3|=-(-3)=3$. So $|x|\ge0$ for every real number $x$.

Absolute values behave well under multiplication: for all real numbers $c$ and $x$,
\[
|cx|=|c|\,|x|.
\]
To check this rule, we go through the possible signs of $c$ and $x$.
\begin{itemize}
\item If $c=0$ or $x=0$, then $cx=0$, and both sides are zero.
\item If $c>0$ and $x>0$, then $cx>0$, and both sides equal $cx$.
\item If $c<0$ and $x<0$, then $cx>0$, so $|cx|=cx$, while $|c|\,|x|=(-c)(-x)=cx$.
\item If one of $c$ and $x$ is positive and the other is negative, then $cx<0$, so $|cx|=-cx$. When $c>0$ and $x<0$, we get $|c|\,|x|=c(-x)=-cx$; when $c<0$ and $x>0$, we get $|c|\,|x|=(-c)x=-cx$.
\end{itemize}
These cases cover every possibility, so the rule always holds. Taking $c=-1$ gives $|-x|=|-1|\,|x|=|x|$: changing the sign of a number does not change its distance from zero. Later chapters use the product rule often, for instance to see how an error changes when it is multiplied by a constant.

For a number $b\ge0$, the statement $|x|\le b$ says that $x$ is at distance at most $b$ from zero, that is,
\[
|x|\le b\quad\Longleftrightarrow\quad -b\le x\le b.
\]
To check this, consider the two cases in the formula for $|x|$. If $x\ge0$, then $|x|=x$, so $|x|\le b$ says $x\le b$; and $-b\le x$ holds automatically, since $-b\le0\le x$. If $x<0$, then $|x|=-x$, so $|x|\le b$ says $-x\le b$, which becomes $x\ge-b$ after multiplying both sides by $-1$; and $x\le b$ holds automatically, since $x<0\le b$.

For example, $|x|\le2$ means $-2\le x\le2$, which is the conjunction ``$-2\le x$ and $x\le2$.'' By De Morgan's law, its negation is a disjunction. The negation of $-2\le x$ is $x<-2$, and the negation of $x\le2$ is $x>2$, so the negation of $|x|\le2$ is ``$x<-2$ or $x>2$.'' In other words, $|x|>2$ means that $x<-2$ or $x>2$. This condition has two separate pieces, one on each side of zero. This two-sided form is exactly what appears when an inequality involves a square, as we now show.

First, squaring preserves order for nonnegative numbers. Suppose $0\le a<b$. Multiplying out one bracket at a time gives
\[
(b-a)(b+a)=b^2+ba-ab-a^2=b^2-a^2.
\]
Both factors on the left are positive: $b-a>0$ because $a<b$, and $b+a>0$ because $b$ is positive (it is larger than $a\ge0$) and $a$ is nonnegative. A product of two positive numbers is positive, so $b^2-a^2>0$, that is, $a^2<b^2$. In the same way, if $0\le a\le b$, then both factors are nonnegative, so $b^2-a^2\ge0$ and $a^2\le b^2$. The condition that the numbers are nonnegative cannot be dropped: $-3<2$, but $(-3)^2=9$ is larger than $2^2=4$.

Second, $x^2=|x|^2$ for every real number $x$. If $x\ge0$, this says $x^2=x^2$. If $x<0$, it says $x^2=(-x)^2$, which is true because $(-x)(-x)=x\cdot x$.

Now we can solve the inequality $x^2>4$ completely. Since $4=2^2$, we compare $|x|$ with $2$.
\begin{itemize}
\item If $|x|>2$, apply the strict rule with $a=2$ and $b=|x|$: we get $4<|x|^2=x^2$.
\item If $|x|\le2$, apply the non-strict rule with $a=|x|$ and $b=2$: we get $x^2=|x|^2\le4$.
\end{itemize}
The first case shows that $|x|>2$ implies $x^2>4$. The second case shows that $|x|\le2$ implies $x^2\le4$, which is the contrapositive of ``$x^2>4$ implies $|x|>2$.'' Together, the two cases prove
\[
x^2>4\quad\Longleftrightarrow\quad|x|>2\quad\Longleftrightarrow\quad x<-2\ \text{or}\ x>2.
\]
For instance, $x=-3$ satisfies $x^2=9>4$ even though $x$ is not greater than $2$, while $x=1.5$ gives $x^2=2.25$, which is not greater than $4$. A tempting shortcut, ``take square roots to get $x>2$,'' loses the whole negative half of the answer.

Ranges of real numbers like these have a standard notation. For real numbers $a<b$, an \emph{interval} is one of the following sets:
\begin{itemize}
\item $[a,b]$ is the set of all real $x$ with $a\le x\le b$;
\item $(a,b)$ is the set of all real $x$ with $a<x<b$;
\item $[a,b)$ is the set of all real $x$ with $a\le x<b$, and $(a,b]$ is the set of all real $x$ with $a<x\le b$.
\end{itemize}
A square bracket includes its endpoint, and a round bracket excludes it. So $1$ and $3$ belong to $[1,3]$ but not to $(1,3)$, while $2$ belongs to both. An interval may also have no endpoint on one side. The notation $[0,\infty)$ means all real $x$ with $x\ge0$, and $(-\infty,-2)$ means all real $x$ with $x<-2$. The symbol $\infty$, read ``infinity,'' only says that there is no endpoint on that side. Infinity is not a real number, so it is never placed next to a square bracket. In this notation, $|x|\le2$ says that $x$ belongs to $[-2,2]$, and $x^2>4$ says that $x$ belongs to $(-\infty,-2)$ or to $(2,\infty)$.

Round brackets are also used for an \emph{ordered pair}, such as $(1,3)$, which lists two entries in a definite order; Chapter~\ref{ch:maps} uses ordered pairs. The surrounding sentence tells you whether $(1,3)$ means an interval or a pair.

\begin{principle}{Habit 2: keep a statement contract}
Before working on a statement, write down its \emph{contract}: the domain of each variable, the hypotheses, the conclusion, and the order in which the quantities are chosen. Whenever the statement changes, for example when you or an assistant repair it or generalize it, compare the new version with the contract. Which hypothesis became weaker or stronger? Did the conclusion become weaker or stronger? A changed statement can be true and still fail to answer the original question. In Chapter~\ref{ch:start}, an assistant that replaced ``integers'' by ``positive integers'' gave a correct proof of a weaker statement, one that says nothing about a sum such as $-3+5$.
\end{principle}
